% ============================================================
\section{Utility Tree y priorización}\label{sec:utility}
% ============================================================

El Utility Tree (árbol de utilidad) organiza las necesidades de calidad del sistema en cuatro niveles: la utilidad del juego como raíz, las características de calidad de la sección~\ref{sec:calidad}, el refinamiento de cada característica en la subcaracterística que realmente importa y, en las hojas, los escenarios de la sección~\ref{sec:escenarios}. Cada hoja recibe dos valoraciones, importancia y sensibilidad arquitectónica, que sirven para decidir qué escenarios deben guiar la arquitectura y cuáles pueden resolverse después, en el diseño detallado.

La prioridad se asigna a cada escenario y no a cada característica, porque una misma característica agrupa escenarios con consecuencias muy distintas. Por ejemplo, QA-05 incluye ESC-12, que decide si cliente y servidor comparten el motor de reglas, y ESC-15, que se resuelve sacando la marca y el arte del código.

% ------------------------------------------------------------
\subsection{Escalas de valoración}

Cada escenario se valora con un par (I, S): \textbf{I} es la importancia para los stakeholders y \textbf{S} la sensibilidad arquitectónica. Ambas usan tres niveles: A (alta), M (media) y B (baja).

\textbf{Importancia.} Mide qué tan grave es para los stakeholders que el escenario no se cumpla. La deciden los intereses, restricciones y requisitos impuestos de las secciones~\ref{sec:stakeholders} a~\ref{sec:preguntas}, no el arquitecto.

\begin{tabla}{|L{1.4cm}|Y|}
\hline
\filacab \cab{Nivel} & \cab{Criterio} \\ \hline
\endhead
A & Si no se cumple, se incumple una restricción (CON) o un requisito impuesto por el cliente (REQ), o un stakeholder con poder de decisión rechaza el producto: el cliente, el padre o tutor, el asesor legal o el titular de la propiedad intelectual. \\ \hline
M & Si no se cumple, el producto sigue siendo aceptable, pero algún stakeholder sufre una pérdida notable, o el escenario depende de una pregunta abierta con una probabilidad real de resolverse a favor. \\ \hline
B & El escenario depende de una decisión que puede no tomarse nunca y, si no se cumple, el daño no afecta a la entrega ni a la seguridad de los usuarios. \\ \hline
\end{tabla}

\textbf{Sensibilidad arquitectónica.} Mide cuánto depende el cumplimiento del escenario de las decisiones estructurales: dónde vive la autoridad y el estado de la partida, cómo se parten los componentes, cómo se comunican y cuándo se persiste la información. La decide el arquitecto.

\begin{tabla}{|L{1.4cm}|Y|}
\hline
\filacab \cab{Nivel} & \cab{Criterio} \\ \hline
\endhead
A & Hay al menos dos alternativas estructurales razonables, y una lo cumple y otra no, o lo cumplen de forma muy distinta. Cambiar de alternativa después obliga a modificar varios componentes. \\ \hline
M & Exige una decisión de diseño acotada: una interfaz, un estado adicional o una política dentro de un componente, o una convención que todas las capas deben respetar. Puede introducirse más tarde con un costo moderado sin rehacer la estructura. \\ \hline
B & Se cumple con implementación, contenido o configuración dentro de cualquier arquitectura razonable. No ayuda a elegir entre alternativas. \\ \hline
\end{tabla}

\textbf{Por qué sensibilidad y no dificultad.} Que un escenario sea difícil no significa que la arquitectura importe para cumplirlo. Detectar insultos disfrazados en ESC-04 es difícil, pero la dificultad está en el algoritmo del filtro y ninguna estructura la resuelve. Al revés, ESC-14 no es difícil hoy: basta con no fijar en el código dónde está el servidor. Pero si esa decisión no se toma ahora, agregar la LAN después obliga a separar de nuevo servidor, configuración y red. La dificultad mide el esfuerzo de construir; la sensibilidad mide cuánto condiciona el escenario a la estructura, que es lo que interesa para priorizar la arquitectura.

% ------------------------------------------------------------
% (el título empieza en la página del diagrama si no cabe junto con él)
\Needspace*{0.66\textheight}
\subsection{Árbol de utilidad}

El diagrama muestra el árbol completo. Un escenario que pone a prueba dos características aparece bajo ambas con la misma valoración. Se resaltan los escenarios (A, A), que guían la elección de la arquitectura, y el único escenario (M, A), que exige una decisión temprana aunque su importancia sea media.

\begin{center}
  % ============================================================
%  Utility Tree: utilidad -> característica -> refinamiento -> escenario (I, S)
%  Requiere \usepackage[edges]{forest} y los colores del preámbulo.
% ============================================================
\definecolor{lineaArbol}{HTML}{8497B0}
\definecolor{fondoAA}{HTML}{C9D6EE}
\definecolor{bordeMA}{HTML}{C55A11}
\definecolor{fondoMA}{HTML}{FBE5D6}

\begin{forest}
  for tree={
    grow'=0, forked edges, fork sep=5pt,
    anchor=west, child anchor=west, parent anchor=east,
    draw=lineaArbol, line width=0.4pt, rounded corners=1.2pt,
    font=\fontsize{6.4}{7.6}\selectfont,
    node options={align=left}, inner xsep=3pt, inner ysep=2.2pt,
    s sep=2.2pt, l sep=9pt,
    edge={draw=lineaArbol, line width=0.4pt},
  },
  where level=0{fill=titulo, text=white, font=\fontsize{7.5}{9}\selectfont\bfseries,
                text width=1.35cm, node options={align=center}}{},
  where level=1{fill=encabezado, font=\fontsize{6.4}{7.6}\selectfont\bfseries, text width=2.25cm}{},
  where level=2{fill=white, text width=2.85cm}{},
  where level=3{fill=white, text width=7.3cm}{},
  % Resaltado de las hojas según su prioridad
  aa/.style={fill=fondoAA, draw=titulo, line width=0.7pt},
  ma/.style={fill=fondoMA, draw=bordeMA, line width=0.7pt},
  [Utilidad del juego
    [QA-01 Desempeño
      [Comportamiento temporal
        [{ESC-01 Colocar una barrera con el tablero avanzado \hfill\textbf{(A, A)}}, aa]]]
    [QA-02 Usabilidad
      [Protección contra errores de usuario
        [{ESC-02 Un niño intenta una barrera que encierra a su rival \hfill(A, B)}]]
      [Capacidad de aprendizaje
        [{ESC-03 Un niño inicia sesión en un equipo nuevo con segundo factor \hfill(A, M)}]]]
    [QA-03 Seguridad
      [Integridad
        [{ESC-05 Un cliente alterado intenta hacer trampa \hfill\textbf{(A, A)}}, aa]]
      [Confidencialidad
        [{ESC-06 Alguien captura el tráfico en una red compartida \hfill\textbf{(A, A)}}, aa]]
      [Autenticidad
        [{ESC-03 Un niño inicia sesión en un equipo nuevo con segundo factor \hfill(A, M)}]]]
    [QA-04 Localizabilidad
      [Adaptabilidad a otros idiomas y regiones
        [{ESC-11 Se agrega un tercer idioma con otro sistema de escritura \hfill(A, M)}]]]
    [QA-05 Mantenibilidad
      [Modificabilidad
        [{ESC-12 El diseñador cambia los parámetros de las reglas \hfill\textbf{(A, A)}}, aa]
        [{ESC-14 Se aprueba el modo LAN con el modo en línea ya entregado \hfill\textbf{(M, A)}}, ma]
        [{ESC-15 Hay que cambiar nombre, arte y recursos por motivos legales \hfill(A, B)}]
        [{ESC-16 Se cambia el mecanismo del segundo factor \hfill(M, M)}]]
      [Modularidad
        [{ESC-13 Se recorta la IA en la semana 10 \hfill(A, M)}]]]
    [QA-06 Fiabilidad
      [Recuperabilidad y tolerancia a fallos de red
        [{ESC-08 Se corta la conexión del jugador a mitad de su turno \hfill\textbf{(A, A)}}, aa]]
      [Recuperabilidad y disponibilidad de datos
        [{ESC-09 SQL Server se cae con partidas en curso \hfill\textbf{(A, A)}}, aa]]
      [Tolerancia a fallos de terceros
        [{ESC-10 El servicio de correo deja de responder durante registros \hfill(M, M)}]]]
    [QA-07 Capacidad de prueba
      [Reproducción sin interfaz
        [{ESC-12 El diseñador cambia los parámetros de las reglas \hfill\textbf{(A, A)}}, aa]
        [{ESC-17 Regresión del incremento semanal sin interfaz \hfill(A, M)}]]]
    [QA-08 Responsabilidad
      [Registro de acciones y sanciones
        [{ESC-09 SQL Server se cae con partidas en curso \hfill\textbf{(A, A)}}, aa]
        [{ESC-07 Varios jugadores reportan en grupo a un inocente \hfill(M, M)}]]]
    [QA-09 Protección del menor
      [Operación a prueba de fallos
        [{ESC-04 Llega un mensaje inapropiado al chat de una partida \hfill(A, M)}]]
      [Identificación de riesgos
        [{ESC-07 Varios jugadores reportan en grupo a un inocente \hfill(M, M)}]]]
    [QA-10 Instalabilidad
      [Instalación en equipos limpios
        [{ESC-18 La plataforma de distribución revisa el paquete \hfill(B, B)}]]]
  ]
\end{forest}

\vspace{4pt}
{\fontsize{6.8}{8}\selectfont
\tikz[baseline=-0.6ex]\node[draw=titulo, line width=0.7pt, fill=fondoAA, rounded corners=1.2pt, minimum width=4mm, minimum height=2.6mm]{};~(A, A): guían la elección de la arquitectura
\quad
\tikz[baseline=-0.6ex]\node[draw=bordeMA, line width=0.7pt, fill=fondoMA, rounded corners=1.2pt, minimum width=4mm, minimum height=2.6mm]{};~(M, A): decisión temprana
\quad
(I, S) = (importancia, sensibilidad arquitectónica)}

\end{center}

La tabla recorre el mismo árbol de la raíz a las hojas, en un formato más fácil de consultar.

\begin{tabla}{|L{3.3cm}|L{4.2cm}|Y|L{1.2cm}|}
\hline
\filacab \cab{Característica} & \cab{Refinamiento} & \cab{Escenario} & \cab{(I, S)} \\ \hline
\endhead
QA-01 Desempeño & Comportamiento temporal & ESC-01 Colocar una barrera con el tablero avanzado & (A, A) \\ \hline
QA-02 Usabilidad & Protección contra errores de usuario & ESC-02 Un niño intenta una barrera que encierra a su rival & (A, B) \\ \cline{2-4}
 & Capacidad de aprendizaje & ESC-03 Un niño inicia sesión en un equipo nuevo con segundo factor & (A, M) \\ \hline
QA-03 Seguridad & Integridad & ESC-05 Un cliente alterado intenta hacer trampa & (A, A) \\ \cline{2-4}
 & Confidencialidad & ESC-06 Alguien captura el tráfico en una red compartida & (A, A) \\ \cline{2-4}
 & Autenticidad & ESC-03 Un niño inicia sesión en un equipo nuevo con segundo factor & (A, M) \\ \hline
QA-04 Localizabilidad & Adaptabilidad a otros idiomas y regiones & ESC-11 Se agrega un tercer idioma con otro sistema de escritura & (A, M) \\ \hline
QA-05 Mantenibilidad & Modificabilidad & ESC-12 El diseñador cambia los parámetros de las reglas & (A, A) \\ \cline{3-4}
 & & ESC-14 Se aprueba el modo LAN con el modo en línea ya entregado & (M, A) \\ \cline{3-4}
 & & ESC-15 Hay que cambiar nombre, arte y recursos por motivos legales & (A, B) \\ \cline{3-4}
 & & ESC-16 Se cambia el mecanismo del segundo factor & (M, M) \\ \cline{2-4}
 & Modularidad & ESC-13 Se recorta la IA en la semana 10 & (A, M) \\ \hline
QA-06 Fiabilidad & Recuperabilidad y tolerancia a fallos de red & ESC-08 Se corta la conexión del jugador a mitad de su turno & (A, A) \\ \cline{2-4}
 & Recuperabilidad y disponibilidad de datos & ESC-09 SQL Server se cae con partidas en curso & (A, A) \\ \cline{2-4}
 & Tolerancia a fallos de terceros & ESC-10 El servicio de correo deja de responder durante registros & (M, M) \\ \hline
QA-07 Capacidad de prueba & Reproducción sin interfaz & ESC-12 El diseñador cambia los parámetros de las reglas & (A, A) \\ \cline{3-4}
 & & ESC-17 Regresión del incremento semanal sin interfaz & (A, M) \\ \hline
QA-08 Responsabilidad & Registro de acciones y sanciones & ESC-09 SQL Server se cae con partidas en curso & (A, A) \\ \cline{3-4}
 & & ESC-07 Varios jugadores reportan en grupo a un inocente & (M, M) \\ \hline
QA-09 Protección del menor & Operación a prueba de fallos & ESC-04 Llega un mensaje inapropiado al chat de una partida & (A, M) \\ \cline{2-4}
 & Identificación de riesgos & ESC-07 Varios jugadores reportan en grupo a un inocente & (M, M) \\ \hline
QA-10 Instalabilidad & Instalación en equipos limpios & ESC-18 La plataforma de distribución revisa el paquete & (B, B) \\ \hline
\end{tabla}

% ------------------------------------------------------------
\subsection{Justificación de las valoraciones}

\begin{tabla}{|L{1.25cm}|L{1.05cm}|L{5.6cm}|Y|}
\hline
\filacab \cab{ID} & \cab{(I, S)} & \cab{Por qué esa importancia} & \cab{Por qué esa sensibilidad} \\ \hline
\endhead
ESC-01 & (A, A) &
REQ-02 fija el límite de 1 segundo por imposición del cliente, y colocar una barrera es una de las acciones más frecuentes del juego (STK-01, CRN-02). &
El segundo se reparte entre tres decisiones estructurales: si valida solo el servidor o también el cliente (TEN-01), si la jugada se persiste antes o después de responder (TEN-06) y cómo viaja el mensaje por el protocolo propio. Cambiar cualquiera de ellas afecta al cliente, al servidor y a la persistencia. \\ \hline
ESC-02 & (A, B) &
REQ-07 exige que un niño de 8 años entienda el juego, y esta es la jugada inválida más frecuente: si no la entiende, abandona (CRN-01). &
Se resuelve con diseño de interfaz y con que el motor de reglas informe qué camino quedaría cerrado, algo que cualquier arquitectura con un motor de reglas puede ofrecer. La medida de 5 segundos se cumple incluso validando solo en el servidor, así que no favorece a ninguna alternativa. \\ \hline
ESC-03 & (A, M) &
Junta dos restricciones no negociables, REQ-03 y REQ-07. Incumplir cualquiera de las dos invalida la entrega. &
Afecta al módulo de cuentas: estados de sesión adicionales, aprobación del tutor y el mecanismo del segundo factor detrás de una interfaz (CRN-26). No cambia cómo se reparten cliente y servidor ni el protocolo de partida. \\ \hline
ESC-04 & (A, M) &
El filtro es obligatorio (REQ-04), y el padre o tutor puede vetar el juego si su hijo lee un insulto (CRN-04). &
Que el filtro se aplique en el servidor se deriva de la autoridad del servidor, que ya decide ESC-05. Lo que agrega este escenario, suspender el chat cuando el filtro falla, es una política dentro del componente de chat. \\ \hline
ESC-05 & (A, A) &
El modo en línea es obligatorio (REQ-01), y un cliente alterado que hace trampa rompe cualquier partida. STK-15 da por hecho que ocurrirá (CRN-23). &
Pone a prueba la decisión más cara de cambiar: quién tiene la autoridad sobre el estado de la partida (CRN-10). Depende además del diseño del protocolo propio (CON-03): delimitado, versionado y validación de mensajes. Pasar de un cliente autoritativo a un servidor autoritativo obliga a rehacer cliente, servidor y protocolo. \\ \hline
ESC-06 & (A, A) &
Sin confidencialidad, el segundo factor de REQ-03 no protege nada y quedan expuestos datos de menores (CRN-05, CRN-22). &
Obliga a elegir entre cifrar la conexión completa o cifrar cada mensaje dentro del protocolo propio (Q-09), y a diseñar sesiones que no se puedan reutilizar. La elección cambia el formato de los mensajes en ambos extremos y consume parte del segundo de ESC-01. \\ \hline
ESC-07 & (M, M) &
El daño a un inocente es real (TEN-02), pero depende de cómo se resuelva Q-06 y no afecta a las partidas ni a la entrega. &
Exige registrar reportes y sanciones con su contexto y una política de retención. Si Q-06 incluye revisión humana, agrega un módulo administrativo, pero separado del núcleo de partida. \\ \hline
ESC-08 & (A, A) &
Es la falla más común, porque los niños juegan desde redes domésticas, y es la situación que describe CRN-06. &
Depende de dónde vive el estado autoritativo de la partida y de si el protocolo propio incluye detección de actividad, porque TCP no avisa de las conexiones medio abiertas (CON-17). Ambas decisiones atraviesan cliente, servidor y protocolo. \\ \hline
ESC-09 & (A, A) &
Perder cuentas o progreso ya confirmados es inaceptable para el jugador y para STK-13, y CRN-19 exige definir de antemano qué se puede perder. &
Depende de cuándo se persiste cada jugada (TEN-06): escritura síncrona, asíncrona o con un registro intermedio. Cada alternativa cambia el flujo entre el servidor de partidas y la base de datos, y condiciona también a ESC-01. \\ \hline
ESC-10 & (M, M) &
Solo aplica si Q-03 o Q-04 meten el correo en el registro o en el acceso, e incluso entonces afecta a quien se registra, no a las partidas. &
Obliga a tratar el registro como un proceso con estado pendiente y reintentos, pero dentro del módulo de cuentas. \\ \hline
ESC-11 & (A, M) &
El cambio de idioma es obligatorio (REQ-06) y la distribución internacional es un objetivo del proyecto. &
No cambia la partición en componentes, pero impone convenciones que todas las capas deben respetar: catálogos de textos, Unicode en el protocolo y en SQL Server (CON-05) y contenedores flexibles en la interfaz. Es barata al inicio y cara al final, pero no distingue entre alternativas estructurales. \\ \hline
ESC-12 & (A, A) &
El diseñador cambiará los parámetros decenas de veces (CRN-13), y Q-08 puede obligar a introducir variantes por motivos legales (CON-15). &
Obliga a decidir si cliente y servidor comparten un único motor de reglas o duplican la lógica (CRN-12, TEN-01), y a sacar los parámetros del código. Si se duplica el motor y después se descubre que los dos lados no coinciden, hay que reescribir ambos. \\ \hline
ESC-13 & (A, M) &
Con un plazo de 14 semanas (CON-13), no poder recortar la IA sin dañar lo obligatorio pone en riesgo la entrega (CRN-07). &
Basta con abstraer el concepto de oponente detrás de una interfaz (TEN-04). Es una decisión acotada al punto donde se decide contra quién se juega. \\ \hline
ESC-14 & (M, A) &
La LAN no está decidida (Q-02), y el modo en línea, que sí es obligatorio, funciona sin ella. &
Depende de la topología de despliegue: si el servidor de partidas puede ejecutarse fuera de la infraestructura remota y si el extremo de red se configura desde fuera (CRN-20). Si se asume un servidor siempre remoto, agregar la LAN obliga a separar de nuevo servidor, configuración y red. \\ \hline
ESC-15 & (A, B) &
El asesor legal y STK-18 pueden exigir el cambio o detener el proyecto (CON-15), y el nombre todavía no está decidido (Q-07). &
Se cumple con el mismo mecanismo de recursos externos que exige ESC-11: marca, arte y textos fuera del código. No requiere ninguna decisión estructural adicional. \\ \hline
ESC-16 & (M, M) &
Solo ocurre si el primer mecanismo no cumple ESC-03, y afecta al acceso, no a las partidas. &
Requiere que la verificación esté detrás de una interfaz (CRN-26) y una migración de cuentas, ambas dentro del módulo de cuentas. \\ \hline
ESC-17 & (A, M) &
Sin una regresión rápida no es viable entregar un incremento auditable cada semana (CON-14), y su resultado es la evidencia del auditor (CRN-08). &
La separación entre motor e interfaz ya la exige ESC-12. Lo que agrega este escenario, que las partidas sean deterministas y queden registradas jugada por jugada, se resuelve dentro del motor y del registro. \\ \hline
ESC-18 & (B, B) &
Solo aplica si se decide publicar (Q-10), y nada indica que vaya a ocurrir dentro del plazo. &
Se resuelve con la configuración del empaquetado, distribuyendo .NET junto con el juego, sin afectar a la estructura. \\ \hline
\end{tabla}

% ------------------------------------------------------------
\subsection{Casos en que importancia y sensibilidad no coinciden}

Importancia y sensibilidad responden a preguntas distintas, y en varios escenarios dan valores opuestos. Estos casos son los que muestran por qué no basta con una sola escala.

\begin{itemize}
  \item \textbf{Alta importancia, baja sensibilidad: ESC-02 y ESC-15.} Son de los escenarios que más importan: un niño que no entiende la regla abandona el juego, y un problema legal puede detener el proyecto. Sin embargo, ninguna alternativa de arquitectura los cumple mejor que otra. Se resuelven con diseño de interfaz y con recursos fuera del código. Siguen siendo criterios de aceptación obligatorios, pero no sirven para elegir la arquitectura.
  \item \textbf{Alta importancia, sensibilidad media: ESC-03, ESC-04, ESC-11, ESC-13 y ESC-17.} Todos derivan de restricciones obligatorias, pero su impacto queda dentro de un componente o de una convención. Se atienden en el diseño detallado de cada módulo, siempre que la estructura elegida no los impida.
  \item \textbf{Importancia media, alta sensibilidad: ESC-14.} La LAN no es obligatoria y el juego se puede entregar sin ella. Aun así, que se pueda agregar depende de una decisión estructural que es barata hoy y cara después. Por eso recibe atención temprana, en coherencia con TEN-04, que eligió el extremo de red como uno de los dos puntos que conviene abstraer.
\end{itemize}

% ------------------------------------------------------------
\subsection{Escenarios que reciben mayor atención}

\begin{tabla}{|L{3.1cm}|L{3.3cm}|Y|}
\hline
\filacab \cab{Grupo} & \cab{Escenarios} & \cab{Tratamiento} \\ \hline
\endhead
(A, A) & ESC-01, ESC-05, ESC-06, ESC-08, ESC-09, ESC-12 &
Guían la elección de la arquitectura. Una alternativa que no pueda cumplir alguno se descarta, aunque sea mejor en otros aspectos. Son los candidatos a requisitos arquitectónicamente significativos. \\ \hline
(M, A) & ESC-14 &
No guía la elección, pero la decisión que lo mantiene posible se toma desde el inicio, porque postergarla es lo que lo vuelve caro. \\ \hline
(A, M) y (A, B) & ESC-02, ESC-03, ESC-04, ESC-11, ESC-13, ESC-15, ESC-17 &
Criterios de aceptación obligatorios que se resuelven en el diseño detallado de cada componente, dentro de la estructura elegida. \\ \hline
(M, M) y (B, B) & ESC-07, ESC-10, ESC-16, ESC-18 &
Se atienden cuando se resuelva la pregunta abierta de la que dependen: Q-06, Q-03, Q-04 o Q-10. \\ \hline
\end{tabla}

Los seis escenarios (A, A) no se pueden evaluar por separado, porque compiten entre sí. ESC-01 pide responder en menos de un segundo, mientras que ESC-05 y ESC-06 agregan validación y cifrado a ese mismo segundo (TEN-01), y ESC-09 pide persistir lo suficiente para no perder datos, lo que también le cuesta tiempo (TEN-06). ESC-08 y ESC-12 dependen de la misma decisión que ESC-05: dónde vive el estado autoritativo y si el motor de reglas es uno solo para cliente y servidor. Por eso, cualquier alternativa de arquitectura se evalúa contra los seis a la vez, y no se da por buena una solución para uno si rompe otro.
